\section{Cubic gaps and the role of unequal means}\label{sec:cubic}
The divergence in Theorem~\ref{thm:positive-growth} requires increasing
degree. Degree three already distinguishes positive multilinear functions
from the positive bilinear factor-two theorem, however. We prove
\begin{equation}\label{eq:cubic-interval}
 \frac{1610000}{743033}\le R_3\le\frac{31}{12}.
\end{equation}
The lower endpoint is a bound on a supremum supplied by a sequence of
certificates, not a claimed attained finite optimum. We also give exact
finite hull values and a homogeneous unit-coefficient example with only
52 variables. The exact value of $R_3$ remains open.

\subsection{Endpoint orientation and a cubic mixture}
One particularly simple global law applies in every degree. Draw
$U\in[0,1]$ uniformly and independently orient each coordinate left or
right with equal probability. A left coordinate is the indicator of
$[0,x_i]$ and a right coordinate of $[1-x_i,1]$. Call this law $O$.
Every coordinate has mean $x_i$.

\begin{proposition}[Endpoint-orientation bound]\label{prop:endpoint-orientation}
For maximum degree $d\ge2$, on every nonnegative box,
\begin{equation}\label{eq:endpoint-bound}
 T\le\frac{d-1}{1-2^{-(d-1)}}H.
\end{equation}
In particular the constant is $8/3$ for cubics and two for bilinear terms.
\end{proposition}
\begin{proof}
On the unit cube, fix a degree-$(k+1)$ term with minimum mean $u$ at its
anchor. Conditional on the anchor orientation, a same-oriented coordinate
imposes no further restriction on its success interval. An opposite
coordinate excludes an interval of length $a_j=\min(u,1-x_j)$ at the
opposite end. These excluded intervals are nested. If
$a_{(1)}\ge\cdots\ge a_{(k)}$, the first selected index has probability
$2^{-j}$ of being $j$, and hence
\[
 \E_O D_e=\sum_{j=1}^k2^{-j}a_{(j)}
 \ge\frac{1-2^{-k}}{k}\sum_ja_j
 \ge\frac{1-2^{-k}}{k}\min(u,\sum_j(1-x_j)).
\]
The first inequality follows by summing
$(2^{-i}-2^{-j})(a_{(i)}-a_{(j)})\ge0$ over pairs, and the second by
$\sum_j\min(u,p_j)\ge\min(u,\sum_jp_j)$. The fraction
$(1-2^{-k})/k$ decreases with $k$, since it is the average of the first
$k$ entries of $1/2,1/4,\ldots$. Sum the simultaneous guarantees and use
Proposition~\ref{prop:box-transfer}.
\end{proof}

\begin{theorem}[Three-law cubic bound]\label{thm:cubic-upper}
Every positive multilinear polynomial of degree at most three on a finite
nonnegative box satisfies $T\le(31/12)H$. On the unit cube, the global
law
\begin{equation}\label{eq:cubic-mixture}
 P=\frac{18}{31}O+\frac6{31}I+\frac7{31}B
\end{equation}
guarantees deficiency at least $12t_e/31$ for every quadratic or cubic
term. Here $I$ is independent rounding. For $B$, draw one uniform $U$,
set each low coordinate to $\ind\{U\le x_i\}$, and, independently
conditional on $U$, let each high coordinate fail with probability $1/2$
when $U\le2(1-x_i)$ and never otherwise. The weights are optimal among
fixed mixtures of these three laws for a uniform termwise guarantee.
\end{theorem}
\begin{proof}
The definition of $B$ gives each high coordinate failure mean $1-x_i$;
thus all three laws have the prescribed means. Consider sorted cubic
means $u\le v\le w$, and set $t=\min(u,2-v-w)$. Terms with $t=0$
need no estimate. Write $D_O,D_I,D_B$ for their expected deficiencies.

If $u\le v\le1/2$, then $t=u$. Opposite anchor and second-coordinate
orientations give $D_O\ge u/2$, while $D_I=u(1-vw)\ge u/2$.
These two contributions alone give $12t/31$.

If exactly one coordinate is low, set $a=1-v\ge b=1-w$. Directly
integrating the nested exclusion and failure intervals gives
\[
 D_O=\tfrac12\min(u,a)+\tfrac14\min(u,b),\qquad
 D_B=\tfrac12\min(u,2a)+\tfrac14\min(u,2b).
\]
We verify $18D_O+7D_B\ge12\min(u,a+b)$ in all cases. If $a\le u/2$,
the left side is $16a+8b\ge12(a+b)$. If $b\ge u/2$, it is at least
$18(3u/8)+7(3u/4)=12u$. If $b\le u/2\le a\le u$, it is
$9a+8b+7u/2$. When $a+b\le u$, use
$3a+4b\le4u-a\le7u/2$; when $a+b\ge u$, use
$9a+8b\ge8u+a\ge17u/2$. Finally, if $b\le u/2$ and $a\ge u$,
the expression is $25u/2+8b\ge12u$. These cases exhaust $a\ge b\ge0$,
including their boundaries. The $I$ contribution is nonnegative.

If all coordinates are high, put $c=1-u\ge a=1-v\ge b=1-w$.
Integrating the $B$ union probabilities $7/8,3/4,1/2$ over intervals of
lengths $2b,2(a-b),2(c-a)$ gives total failure probability
$c+a/2+b/4$. Subtracting the anchor failure probability $c$ yields
\[
 D_B=D_O=a/2+b/4\ge3(a+b)/8\ge3t/8.
\]
Independence gives $D_I=u(a+b-ab)\ge7t/16$. Indeed, set $s=a+b\le1$
and use $ab\le s^2/4$. For $s\le u$ the normalized expression is at
least $u(1-s/4)\ge u(1-u/4)\ge7/16$. For $s\ge u$ it is at least
$s-s^2/4\ge u-u^2/4\ge7/16$. Both comparisons use $u>1/2$ and
monotonicity on the unit interval. Thus the mixture fraction is at least
$(25/31)(3/8)+(6/31)(7/16)=12/31$.

For sorted quadratic means $u\le v$, set $t=\min(u,1-v)$; always
$D_O=t/2$. If both means are low, $D_I\ge t/2$. If exactly one is low,
$D_B=\min(u,2(1-v))/2\ge t/2$. If both are high, $D_B=t/2$ and
$D_I=ut\ge t/2$. The resulting mixture fractions are at least
$12/31$, $25/62$, and $1/2$. This proves the simultaneous guarantee.
Summation and positive expansion give the claimed box bound.

To check optimality in the stated family, let the arbitrary weights of
$(O,I,B)$ be $(w,r,v)$ with $w+r+v=1$. The following test means give the
listed exact or limiting normalized deficiencies:
\[
\begin{array}{c|ccc}
\text{means}&D_O/t&D_I/t&D_B/t\\\hline
(u,1/2),\ 0<u\le1/2&1/2&1/2&0\\
(\varepsilon,1-\varepsilon^2,1-\varepsilon^2),\ \varepsilon\downarrow0
 &3/8&0&3/4\\
(u,1-u/2,1-u/2),\ u\downarrow1/2\text{ from above}&3/8&7/16&3/8
\end{array}
\]
In the second row $t=2\varepsilon^2$ for small $\varepsilon$ and the
independent fraction is $\varepsilon-\varepsilon^3/2$. In the last row
$1/2<u<2/3$ ensures $t=u$ and the independent fraction is $u-u^2/4$.
Therefore a universal fraction $\alpha$ obeys
\[
 \alpha\le\tfrac12-\tfrac v2,\qquad
 \alpha\le\tfrac38+\tfrac{3v}8-\tfrac{3r}8,\qquad
 \alpha\le\tfrac38+\tfrac r{16}.
\]
Average these inequalities with weights $3/31,4/31,24/31$. The variable
coefficients cancel, giving $\alpha\le12/31$.
This conclusion concerns fixed-mixture termwise guarantees; it does not
identify $R_3$, or exclude a stronger coefficient-dependent analysis.
\end{proof}

\subsection{A scalar certificate and the strongest limiting lower bound}
For a set $Q$ of variables, write
$E_k(Q)=\sum_{S\subseteq Q,\,|S|=k}\prod_{i\in S}x_i$.
At a binary vertex this is $\binom{K}{k}$ when $K$ counts successes in
$Q$. Count-polynomial identities below are vertex identities, not
substitutions defining the multilinear functions at continuous points.

\begin{lemma}[Cubic scalar minorant with positive slack]
\label{lem:cubic-scalar}
On $[0,1]^3$, let
\begin{align*}
 F(a,b,c)&=6c^3+27bc^2+18b^2+30ac+20ab+9a^2,\\
 \ell(a,b,c)&=37a+\tfrac{79}{2}b+38c-\tfrac{103}{3}.
\end{align*}
Then $F\ge\ell+\delta_*$, where $\delta_*=901/120000$.
\end{lemma}
\begin{proof}
For fixed $c$, the Hessian of $h=F-\ell$ in $(a,b)$ is
$\left(\begin{smallmatrix}18&20\\20&36\end{smallmatrix}\right)$,
positive definite with determinant 248. For $0\le c\le3/10$, set
$a=1$ and $b_*(c)=13/24-3c^2/4\in(0,1)$. There $\partial_bh=0$ and
$\partial_ah=-49/6+30c-15c^2\le-31/60<0$.
The convex first-order inequality proves that this is the constrained
minimum on the $(a,b)$ square: $(a-1)\partial_ah\ge0$ for $a\le1$.
Its value is
\[
 p(c)=\frac{-972c^4+576c^3+1404c^2-768c+101}{96}.
\]
For $3/10\le c\le1$, minimizing over unrestricted real $a,b$ is a valid
lower bound. Completing the square gives
\[
 r(c)=\frac{-78732c^4+212256c^3-203796c^2+82032c-11275}{2976}.
\]
For example, with $v=30c-37$, $w=27c^2-79/2$ and
$q=6c^3-38c+103/3$, this value is
$q-(18v^2-20vw+9w^2)/248$, which verifies the elimination directly.

Table~\ref{tab:cubic-bernstein} supplies exact degree-four Bernstein
representations. A row with interval $[l,u]$, denominator $D$, and
numerators $v_0,\ldots,v_4$ means
\[
 p(c)\ \text{or}\ r(c)=D^{-1}\sum_{i=0}^4v_i\binom4i
     t^i(1-t)^{4-i},\qquad t=(c-l)/(u-l).
\]
These are rational polynomial identities, verified by expanding in $t$.
The five intervals cover the required ranges. The Bernstein basis is
nonnegative and sums to one, and the minimum of all 25 coefficients
$v_i/D$ is $1802/240000=\delta_*$. Thus both lower polynomials are at
least $\delta_*$ wherever used, proving the claim.
\end{proof}

\begin{table}[htbp]
\centering\small
\begin{tabular}{ccrl}
\toprule
Polynomial&Interval&$D$&$(v_0,v_1,v_2,v_3,v_4)$\\\midrule
$p$&$[0,1/5]$&60000&$(63125,39125,20975,9395,4133)$\\
$p$&$[1/5,3/10]$&240000&$(16532,6008,1802,3788,11597)$\\
$r$&$[3/10,1/2]$&7440000&$(215357,1285415,1434125,1250375,1008125)$\\
$r$&$[1/2,3/4]$&190464&$(25808,18056,7964,9230,15869)$\\
$r$&$[3/4,1]$&190464&$(15869,22508,34520,45920,31040)$\\\bottomrule
\end{tabular}
\caption{Exact scalar positivity certificate for Lemma~\ref{lem:cubic-scalar}.}
\label{tab:cubic-bernstein}
\end{table}

\begin{theorem}[Analytic cubic lower family]\label{thm:cubic-analytic}
For each multiple $m\ge9$ of nine, take three groups $U,V,W$ of $m$
variables and write $A,B,C$ for their sums. Define
\begin{equation}\label{eq:cubic-analytic-family}
 f_m=2E_3(W)+3B E_2(W)+2mE_2(V)
       +\frac{5m}{3}AC+\frac{10m}{9}AB+mE_2(U).
\end{equation}
At group means $1/4,1/2,3/4$, respectively,
\begin{equation}\label{eq:cubic-analytic-finite}
 \frac{T(f_m)}{H(f_m)}\ge
 \frac{161/4-135/(4m)+6/m^2}
 {223/12-\delta_*+135/(4m)+9/m^2}.
\end{equation}
Consequently $R_3\ge1610000/743033$.
\end{theorem}
\begin{proof}
All coefficients in \eqref{eq:cubic-analytic-family} are positive integers.
At binary counts put $a=A/m,b=B/m,c=C/m$. Expansion of binomial
coefficients gives the exact identity
\begin{equation}\label{eq:cubic-count-expansion}
 \frac{18f_m}{m^3}=F(a,b,c)
 -\frac{18c^2+27bc+18b+9a}{m}+\frac{12c}{m^2}
 \ge F(a,b,c)-72/m.
\end{equation}
Every feasible binary law has expected normalized counts
$(1/4,1/2,3/4)$. Lemma~\ref{lem:cubic-scalar} therefore gives
\[
 \frac{18}{m^3}\vex f_m\ge139/6+\delta_*-72/m.
\]
Counting supports and using \eqref{eq:monomial-envelopes} yields
\begin{align}
 \frac{18}{m^3}\cav f_m&=167/4-153/(4m)+9/m^2,
       \label{eq:cubic-analytic-cav}\\
 \frac{18}{m^3}T(f_m)&=161/4-135/(4m)+6/m^2.
       \label{eq:cubic-analytic-tgap}
\end{align}
For clarity, the six scaled upper contributions in the displayed family
order are
\[
 \frac92-\frac{27}{2m}+\frac9{m^2},\quad
 \frac{27}2-\frac{27}{2m},\quad
 9-\frac9m,\quad\frac{15}2,\quad5,\quad
 \frac94-\frac9{4m}.
\]
Only $2E_3(W)$ has a positive termwise lower envelope; its scaled value
is $3/2-9/(2m)+3/m^2$. This proves the two exact formulas. Subtraction
bounds the scaled hull gap by the denominator in
\eqref{eq:cubic-analytic-finite}. It is positive for $m\ge9$, and the
actual hull gap is positive as well: independence gives strictly smaller
expectations than common-threshold rounding for every positive nonlinear
term at interior means. Division is therefore valid. Taking arbitrarily
large multiples of nine yields
\[
 R_3\ge\frac{161/4}{223/12-901/120000}
       =\frac{4830000}{2229099}=\frac{1610000}{743033}.
\]
Dropping the positive slack recovers the weaker $483/223$; the same
certificate proves the stronger reduced fraction above. Neither minorant
is asserted to be attained in expectation, and neither limiting bound
asserts an attained finite gap ratio or convergence of the actual family
ratios to that lower endpoint.
\end{proof}

\subsection{Exact finite cubic witnesses}
Finite examples permit a stronger type of statement: matching primal and
dual certificates determine both envelopes exactly. The three-group
polynomials in Table~\ref{tab:cubic-finite} have group sizes $m$ and
means $(1/4,1/2,3/4)$, and take the form
\begin{equation}\label{eq:cubic-six-family}
 f=c_1E_3(W)+c_2B E_2(W)+c_3E_2(V)+c_4AC+c_5AB+c_6E_2(U).
\end{equation}

\begin{proposition}[Exact finite values]\label{prop:cubic-finite}
The envelope and ratio values in Table~\ref{tab:cubic-finite} are exact.
The 18- and 24-variable polynomials have 212 and 464 monomials,
respectively. Complete rational primal and dual certificates, and a
short executable integer/rational checker, appear in
Appendix~\ref{app:cubic-certificates}.
\end{proposition}
\begin{proof}
At binary counts $A,B,C\in\{0,\ldots,m\}$ the polynomial is
\[
 \Phi_m(A,B,C)=c_1\binom C3+c_2B\binom C2+c_3\binom B2
              +c_4AC+c_5AB+c_6\binom A2.
\]
The appendix certifies an affine minorant at every count state and a
probability distribution of count states with means $(m/4,m/2,3m/4)$
attaining it. Conditional uniform subsets in each group realize every
individual marginal, so these give matching lower-envelope certificates.
The upper and termwise lower values follow directly from the six orbit
sizes and monomial bounds in the appendix. Subtracting yields the ratios.
\end{proof}

\begin{table}[htbp]
\centering\small
\begin{tabular}{rllrrl}
\toprule
$n$&$(c_1,\ldots,c_6)$&$\cav f$&$\sum_e\vex f_e$&$\vex f$&$T/H$\\\midrule
18&$(2,3,9,10,7,7)$&$1647/4$&$10$&$2750/13$&$20891/10411$\\
24&$(2,3,13,12,8,7)$&$971$&$28$&$3572/7$&$6601/3225$\\
192&$(2,3,120,105,70,63)$&$587944$&$20832$&$34172072/105$&$7443345/3445256$\\\bottomrule
\end{tabular}
\caption{Exact cubic witnesses, all with ratio above two.}
\label{tab:cubic-finite}
\end{table}

The 24-variable example also gives a quantitative homogeneous interior
variant. Write it as $f=P+Q$, where $P$ is cubic and $Q$ quadratic,
and put $g(x,z)=P(x)+zQ(x)$ with $z=999/1000$.
All 25 coordinates are interior and all 464 monomials are cubic, with
integer coefficients at most 13. Their upper and lower termwise values
are unchanged, so $\cav g=971$ and $T(g)=943$.
The sum of quadratic coefficients is
\[
 13\binom82+12\cdot8^2+8\cdot8^2+7\binom82=1840.
\]
Thus for any feasible law,
$\E g\ge\vex f-1840/1000$, regardless of dependence involving $Z$.
Consequently
\begin{equation}\label{eq:homogeneous-25}
 H(g)\le3225/7+46/25=80947/175,\qquad
 \frac{T(g)}{H(g)}\ge\frac{165025}{80947}>2.
\end{equation}
This is a certified bound, rather than an exact new hull value.

For completeness, it yields a finite nonconstructive unit-coefficient
example through Theorem~\ref{thm:coefficient-removal}. Normalize by 13,
clone each variable $m=1000$ times, and take error allowance $t=m^3/10$.
The failure probability in \eqref{eq:clone-union-bound} is at most
\begin{equation}\label{eq:finite-25000-probability}
 2(2^{25m}+1)\exp(-m^3/23200)<1.
\end{equation}
Indeed its logarithm is at most $(25m+2)\log2-m^3/23200<0$ at $m=1000$,
even with $\log2<1$. Since $943-2(80947/175)=3131/175$, every
successful sample satisfies
\begin{equation}\label{eq:finite-25000-margin}
 T(G_m)-2H(G_m)\ge m^3\left(\frac{3131}{2275}-\frac12\right)>0.
\end{equation}
It is nonempty and has positive hull gap at its interior point. This
proves existence in 25,000 variables with $464m^3$ candidate monomials;
no successful sample is claimed to be listed. The explicit construction
below is much smaller, while coefficient removal remains useful for
approximating an arbitrary fixed-degree ratio.

\subsection{Two means and an explicit homogeneous unit-coefficient example}
\begin{theorem}[Two-level cubic family]\label{thm:cubic-two-level}
For each positive multiple $m$ of four, take two groups $U,W$ of $m$
variables with means $1/2,3/4$, and define
\begin{equation}\label{eq:two-level-family}
 f_m=A E_2(W)+\frac{5m}{4}E_2(U),\qquad A=\sum_{i\in U}x_i.
\end{equation}
Then
\begin{equation}\label{eq:two-level-ratio}
 \frac{243}{115}(1-1/m)\le\frac{T(f_m)}{H(f_m)}\le\frac{243}{115}.
\end{equation}
For $m=16$ its exact ratio is $135/67$. There is an explicit polynomial
with 52 variables and 4,320 distinct homogeneous cubic monomials, every
coefficient one, at an interior point with ratio at least $2700/1343>2$.
\end{theorem}
\begin{proof}
For real $a,c\in[0,1]$, set $F_2(a,c)=ac^2+5a^2/4$ and
$\ell_2(a,c)=11a/6+20c/27-95/108$. The identity
\begin{equation}\label{eq:two-level-scalar}
 F_2-\ell_2=\frac54\left(a-\frac{11-6c^2}{15}\right)^2
       +\frac{(1-c)(3c-2)^2(3c+7)}{135}
\end{equation}
proves $F_2\ge\ell_2$ on the square. Equality holds at $(1/3,1)$ and
$(5/9,2/3)$; their mixture with weights $1/4,3/4$ has mean
$(1/2,3/4)$ and expected value $16/27$. Thus the scalar envelope value
at that mean is exactly $16/27$.

At binary counts $A=ma,C=mc$,
\[
 \frac{2f_m}{m^3}=F_2(a,c)-\frac{ac+(5/4)a}{m}.
\]
Every feasible law has $\E a=1/2$ and $\E c=3/4$, and $ac\le a$.
It follows that $2\vex f_m/m^3\ge16/27-9/(8m)$. Each monomial
has upper envelope $1/2$ and lower envelope zero; therefore
\begin{equation}\label{eq:two-level-cav}
 \frac{2}{m^3}\cav f_m=\frac{2}{m^3}T(f_m)=\frac98(1-1/m).
\end{equation}
Subtraction gives $2H(f_m)/m^3\le115/216$.
For the reverse bound, choose one of the two equality atoms above, then
sample all coordinates independently conditional on its two probabilities.
Distinct variables in each term have conditional product expectation,
so $2\E[f_m\mid a,c]/m^3=(1-1/m)F_2(a,c)$. This is a feasible law and
gives $2\vex f_m/m^3\le(1-1/m)16/27$. Hence
\[
 \frac{115}{216}(1-1/m)\le\frac{2H(f_m)}{m^3}\le\frac{115}{216}.
\]
Together with \eqref{eq:two-level-cav} this proves
\eqref{eq:two-level-ratio}, including positive denominators and convergence
of the actual ratios to $243/115$. For example, $m=20$ gives the lower
bound $4617/2300>2$ in 40 variables.

At $m=16$, the exact integer-grid certificate is
\begin{equation}\label{eq:two-level-16-dual}
 A\binom C2+20\binom A2\ge214A+\frac{177}{2}C-1686
 \qquad(A,C\in\{0,\ldots,16\}).
\end{equation}
Appendix~\ref{app:cubic-certificates} supplies its integer verification.
Uniform success subsets of sizes $8,12$ give equality, with value 1088.
The concave envelope and termwise gap are 2160, so $H=1072$ and
$T/H=135/67$.

Introduce 20 distinct coordinates $z_k$ and define explicitly
\begin{equation}\label{eq:unit-52}
 g=A E_2(W)+\left(\sum_{k=1}^{20}z_k\right)E_2(U),
 \qquad x_U=1/2,\quad x_W=3/4,\quad z_k=999/1000.
\end{equation}
There are $16\binom{16}{2}+20\binom{16}{2}=4320$ distinct cubic
supports, all with coefficient one. Every term still has upper envelope
$1/2$ and lower envelope zero, so $\cav g=T(g)=2160$. At binary vertices,
\[
 g=f_{16}-\left(\sum_k(1-z_k)\right)E_2(U),\qquad
 0\le E_2(U)\le120.
\]
Every feasible law projects to the original means, giving
$\vex g\ge1088-120(20/1000)=1088-12/5$.
Thus $0<H(g)\le5372/5$ and $T(g)/H(g)\ge2700/1343$.
Positivity of $H$ follows by comparing independent and common-threshold
laws at these interior means.
\end{proof}

The two-level examples show that means in the upper half of the unit
cube do not force factor two. Small upward perturbations of the $1/2$
means preserve a strict ratio above two by continuity, so all means can
be strictly above $1/2$. No quantitative perturbation radius or minimal
dimension claim is implied. The next result explains why at least two
distinct means are needed on the unit cube.

\section{Equal means and classical symmetric envelopes}\label{sec:equal-means}
The envelope of the complete elementary symmetric polynomial is classical.
Sherali's equation~(13) and Theorem~3 \citep[pp.~252--253]{Sherali1997}
give, on the entire unit cube,
\begin{equation}\label{eq:sherali-symmetric}
 \vex E_d(x)=\max\left(0,
  \max_{k=d-1,\ldots,n-1}
  \left\{\binom{k}{d-1}\sum_i x_i-(d-1)\binom{k+1}{d}\right\}\right).
\end{equation}
The following extremal gap formulation is a consequence of that
symmetric-envelope machinery. We give a direct proof at equal means to
make the common-law guarantee explicit.

\begin{theorem}[Exact equal-mean formula]\label{thm:equal-finite}
Fix $n\ge2$ and $0<u<1$. Put $b=\lfloor nu\rfloor$,
$\theta=nu-b$, and, for $2\le d\le n$,
\[
 t_d(u)=\min\{u,(d-1)(1-u)\},\qquad
 q_{n,d}(u)=\frac{(1-\theta)\binom bd+\theta\binom{b+1}d}{\binom nd},
\]
with $\binom{k}{d}=0$ for integers $0\le k<d$. The maximum of $T/H$
over positive multilinear polynomials with a nonzero nonlinear part at
$x=(u,\ldots,u)$ is
\begin{equation}\label{eq:equal-finite}
 R_n^{\mathrm{eq}}(u)=\max_{2\le d\le n}
                \frac{t_d(u)}{u-q_{n,d}(u)}.
\end{equation}
A maximum-degree restriction $2\le D\le n$ restricts this maximum to
$d\le D$. Each candidate ratio is attained by $E_d$.
\end{theorem}
\begin{proof}
Draw $K=b$ or $b+1$ with probabilities $1-\theta,\theta$, then choose
a uniformly random $K$-element success subset. Each mean is $u$ and
each degree-$d$ product expectation is $q_{n,d}(u)$. This single law gives
\[
 H(f;u\boldsymbol1)\ge\sum_ea_e(u-q_{n,|e|}(u)),\qquad
 T(f;u\boldsymbol1)=\sum_ea_et_{|e|}(u).
\]
The quotient is bounded by the largest term ratio, provided its
denominators are positive. To see positivity and attainment, the binary
value of $E_d$ is $\binom Kd$, whose first forward difference is
$\binom K{d-1}$ and is nondecreasing. Its linear interpolation on integer
counts is convex. Jensen's inequality shows that any integer-valued $K$
with mean $nu$ has expected value at least the adjacent-count
interpolation. Our uniform-subset law attains this value with all
individual means correct, establishing $\vex E_d=\binom nd q_{n,d}$.
Its upper envelope is $u\binom nd$. Comparing the minimum with the
independent law, whose product expectation is $u^d$, gives
$q_{n,d}\le u^d<u$. Thus the denominators are positive and every claimed
candidate is attained. The maximum is attained by a maximizing degree.
\end{proof}

\begin{corollary}[Dimension-free equal-mean supremum]\label{cor:equal-infinite}
For $0<u<1$ the exact supremum over dimensions, and the limit of
$R_n^{\mathrm{eq}}(u)$ as $n\to\infty$, is
\begin{equation}\label{eq:equal-infinite}
 R_\infty^{\mathrm{eq}}(u)
 =\max_{k\ge1}\frac{\min\{1,k(1-u)/u\}}{1-u^k}\le2.
\end{equation}
Writing $r_u=u/(1-u)$, a maximizing $k$ is one of
$\max(1,\lfloor r_u\rfloor)$ and $\max(1,\lceil r_u\rceil)$.
For $u\le1/2$ the value is $1/(1-u)$. The uniform constant two is sharp.
\end{corollary}
\begin{proof}
The comparison $q_{n,d}\le u^d$ gives the upper bound by the displayed
maximum after writing $k=d-1$. For fixed $d$,
$q_{n,d}(u)\to u^d$ as $n\to\infty$, since both adjacent counts divided
by $n$ tend to $u$. Thus any fixed-degree candidate is approached by
$E_d$ in increasing dimension.

For $k\ge r_u$ the expression is $1/(1-u^k)$, decreasing with $k$.
For $k\le r_u$ it equals $k/(u+u^2+\cdots+u^k)$, increasing with $k$
because the average of that decreasing sequence decreases. Hence a finite
maximizer occurs at one of the stated two integers. Its fixed-degree
limit proves both the supremum and convergence assertions.

Set $t=k(1-u)/u$. Bernoulli's inequality gives
$u^{-k}=(1+(1-u)/u)^k\ge1+t$, and therefore
\[
 \frac{\min(1,t)}{1-u^k}\le\min(1,t)\frac{1+t}{t}\le2.
\]
When $u\le1/2$ the maximizing integer is $k=1$, giving $1/(1-u)$.
At $u=1/2$, the complete bilinear polynomial for even $n$ has ratio
$2(n-1)/n$, tending to two. This is the classical positive bilinear
sharpness example, now viewed within the equal-mean formula.
\end{proof}

These exact formulas concern equal normalized means on the unit cube.
Positive expansion transfers the degree upper bounds to nonnegative
boxes, but does not identify an original positive-box factorization with
the elementary symmetric unit-cube model. In particular, the precise
finite equal-mean optimization above makes no such additional box claim.
